\documentclass[conference]{IEEEtran}
\IEEEoverridecommandlockouts

\usepackage{cite}
\usepackage{amsmath,amssymb,amsfonts}
\usepackage{graphicx}
\usepackage{textcomp}
\usepackage{xcolor}
\usepackage{booktabs}
\usepackage{array}

\begin{document}

\title{Improved Priority Exchange Server: Critical Evaluation and Transfer\\
{\footnotesize Milestone 3 --- Special Emphasis: Real-Time Systems, HSHL}}

\author{\IEEEauthorblockN{Sumon Boiddo}
\IEEEauthorblockA{\textit{BEng in Electronic Engineering} \\
\textit{Hochschule Hamm-Lippstadt}\\
Matriculation Number: 2220054 \\
Group B, Team 6}}

\maketitle

\begin{abstract}
This document is the Milestone~3 (Critical Evaluation and Transfer) for the
seminar topic \emph{Improved Priority Exchange Server} (IPE). It critically
evaluates IPE: its concrete strengths and limitations, its hidden or
unrealistic assumptions, its run-time and memory scalability, suitable and
unsuitable application scenarios, and proposed extensions and open research
questions. The primary reference is Buttazzo, \textit{Hard Real-Time Computing
Systems}~\cite{buttazzo}, Chapter~6.
\end{abstract}

% ============================================================
\section{Concrete Strengths}
% ============================================================

\begin{itemize}
  \item \textbf{Near-optimal responsiveness, close to EDL.} IPE serves
        aperiodic (emergency) tasks almost as fast as the optimal EDL server;
        it reaches near-optimal response times, matching EDL in almost all
        cases (only rarely is EDL slightly better). It also always runs at the
        highest priority, so a request is served immediately whenever capacity
        is available~\cite{buttazzo}.
  \item \textbf{Low run-time overhead.} The heavy idle-time calculation is done
        offline (the $E$ and $D$ arrays), so at run time IPE only replays the
        stored result. Its run-time cost is low, comparable to the DPE
        server~\cite{buttazzo}.
  \item \textbf{No wasted CPU time.} Through the priority-exchange mechanism and
        spare-time reclaiming, unused server time is lent to periodic tasks and
        reclaimed later, so the processor is used efficiently.
  \item \textbf{Formal schedulability guarantee.} The periodic tasks' hard
        deadlines are always safe as long as the total periodic utilization
        satisfies $U_p \leq 1$; the server automatically uses the remaining
        bandwidth $1 - U_p$. This is an exact EDF condition, not a
        heuristic~\cite{spuri1996}.
\end{itemize}

% ============================================================
\section{Concrete Limitations and Risks}
% ============================================================

\begin{itemize}
  \item \textbf{High memory cost.} Because IPE precomputes everything offline,
        it must store the $E$ and $D$ arrays. When the periods are non-harmonic,
        the hyperperiod (lcm) becomes very large, so the tables become very
        large and the memory use can become prohibitive --- the main risk of
        IPE~\cite{buttazzo}.
  \item \textbf{Requires a known, static task set.} The periodic task set must
        be known in advance to precompute the arrays. If the task set changes
        at run time, the precomputed plan becomes invalid and the guarantees
        are lost, so IPE cannot be used in dynamic systems.
  \item \textbf{Requires EDF.} IPE needs an EDF scheduler, so it cannot run on
        fixed-priority-only (RM) platforms.
  \item \textbf{Implementation complexity.} IPE is more complex than simpler
        servers (such as TBS), because of the priority-exchange bookkeeping and
        the offline precomputation step.
  \item \textbf{Not fully optimal.} IPE is only nearly optimal; in rare cases
        the EDL server is slightly better.
\end{itemize}

% ============================================================
\section{Hidden or Unrealistic Assumptions}
% ============================================================

\begin{itemize}
  \item \textbf{Fixed, known task set.} IPE assumes the periodic task set is
        fully known and fixed offline. Real systems are often dynamic (tasks
        added, removed, or changed at run time), so this assumption may not
        hold.
  \item \textbf{Single processor.} The basic IPE model assumes a single CPU.
        Modern embedded systems are often multi-core, and IPE's basic
        formulation does not extend to them.
  \item \textbf{Accurate WCET.} IPE assumes each task's worst-case execution
        time is known accurately. In reality, WCET is hard to estimate
        precisely, and a task's actual execution time may vary from run to run.
  \item \textbf{Full preemption.} IPE assumes any task can be paused (preempted)
        instantly. Real systems often have non-preemptible sections (critical
        regions, I/O), which the basic model ignores.
  \item \textbf{Implicit deadlines and synchronous release.} IPE assumes
        deadline equal to period ($D_i = T_i$) and that all periodic tasks
        start together at $t = 0$. Real tasks may have different deadlines and
        arbitrary release times.
\end{itemize}

% ============================================================
\section{Scalability Analysis}
% ============================================================

\begin{itemize}
  \item \textbf{Run-time cost --- scales well (low).} At run time IPE only
        replays the precomputed replenishments and performs light
        priority-exchange bookkeeping (comparable to DPE). Because the heavy
        work is done offline, the run-time cost stays low even as the system
        grows.
  \item \textbf{Memory cost --- scales poorly (high).} The $E$ and $D$ tables
        must cover one hyperperiod $H = \mathrm{lcm}(T_1,\dots,T_n)$. As the
        number of tasks grows, or when the periods are non-harmonic, the
        hyperperiod can grow explosively, so the tables and the memory needed
        become very large. Memory is IPE's scalability bottleneck.
  \item \textbf{Offline computation cost.} Computing the EDL idle times offline
        also grows with the hyperperiod, but this is a one-time design-time cost
        and does not affect run-time performance.
  \item \textbf{Overall.} IPE scales well in time but poorly in memory. It is
        best for systems with a small number of harmonically-related tasks
        (small hyperperiod) and does not scale to large or non-harmonic task
        sets.
\end{itemize}

% ============================================================
\section{Suitable Application Scenario}
% ============================================================

\begin{itemize}
  \item A good example is a \textbf{robot (motor) controller}. It runs fixed
        periodic control loops --- for example reading sensors and updating
        motor commands at regular, harmonically-related rates --- which are
        known and fixed at design time and have hard deadlines that must never
        be missed.
  \item Alongside these, it handles occasional soft aperiodic events such as
        operator commands, status requests, or diagnostic messages, which
        should be served quickly but do not have hard deadlines.
  \item IPE fits this scenario well because the periodic task set is fixed and
        known offline (so the $E$ and $D$ arrays can be precomputed); the
        control periods are usually harmonic (so the hyperperiod and tables stay
        small); the controller runs on a single processor with an EDF
        scheduler; and IPE guarantees the periodic control deadlines
        ($U_p \leq 1$) while serving the aperiodic commands with near-optimal
        speed.
\end{itemize}

% ============================================================
\section{Unsuitable Application Scenario}
% ============================================================

\begin{itemize}
  \item A poor example is an \textbf{automotive engine-control unit (ECU)}
        running an AUTOSAR-based operating system, which uses fixed-priority
        scheduling.
  \item IPE fails here because AUTOSAR is fixed-priority (RM-style) and provides
        no EDF scheduler, so IPE cannot run at all; and because, if the set of
        tasks changes at run time, the precomputed $E$ and $D$ arrays become
        invalid and IPE loses its guarantees. On such a platform a
        fixed-priority server (Polling, Deferrable, or Sporadic Server) is the
        appropriate choice instead.
  \item More generally, IPE is unsuitable for any dynamic system (changing task
        set), any multi-core platform, or any system with non-harmonic periods
        on a memory-limited device.
\end{itemize}

% ============================================================
\section{Proposed Extensions and Open Questions}
% ============================================================

\noindent\textbf{Proposed extensions.}
\begin{itemize}
  \item \textbf{Reduced-memory IPE:} compress or bound the $E$ and $D$ tables
        (for example by approximating the EDL idle times), so IPE can be used
        even when the periods are non-harmonic and the hyperperiod is large.
        This addresses IPE's main limitation.
  \item \textbf{Multi-core IPE:} extend the basic single-processor model to
        multiprocessor / multi-core platforms.
  \item \textbf{Adaptive IPE:} allow the task set to change at run time by
        recomputing or incrementally updating the $E$ and $D$ arrays online, so
        IPE can be used in dynamic systems.
  \item \textbf{Robust IPE:} make IPE tolerant to WCET uncertainty, since
        worst-case execution time is hard to estimate exactly in practice.
\end{itemize}

\noindent\textbf{Open questions.}
\begin{itemize}
  \item How can the memory cost of the $E$ and $D$ tables be bounded or reduced
        for realistic non-harmonic task sets?
  \item Can IPE be extended to multi-core while keeping its near-optimal
        responsiveness and its schedulability guarantee?
  \item Can the offline precomputation be made incremental, so the task set can
        change at run time without recomputing everything?
  \item Is there a closed-form worst-case response-time bound for aperiodic
        requests under IPE?
\end{itemize}

% ============================================================
\section{Milestone 3 Checklist Mapping}
% ============================================================

\begin{table}[h]
\centering
\caption{M3 Checklist Mapping}
\renewcommand{\arraystretch}{1.3}
\begin{tabular}{|p{5.5cm}|p{2.0cm}|}
\hline
\textbf{M3 Requirement} & \textbf{Section} \\
\hline
Concrete strengths of IPE & Section 1 \\
\hline
Concrete limitations and risks of IPE & Section 2 \\
\hline
Hidden or unrealistic assumptions in IPE & Section 3 \\
\hline
Runtime, memory, scalability implications & Section 4 \\
\hline
Suitable application example for IPE & Section 5 \\
\hline
Unsuitable application example for IPE & Section 6 \\
\hline
Proposed extensions and open questions & Section 7 \\
\hline
\end{tabular}
\end{table}

% ============================================================
% References
% ============================================================

\begin{thebibliography}{9}

\bibitem{buttazzo}
G.~C.~Buttazzo,
\textit{Hard Real-Time Computing Systems: Predictable Scheduling Algorithms
and Applications}, 3rd~ed.
Springer, 2011, ch.~6 (Dynamic Priority Servers).

\bibitem{spuri1994}
M.~Spuri and G.~Buttazzo,
``Efficient Aperiodic Service under Earliest Deadline Scheduling,''
in \textit{Proc. IEEE Real-Time Systems Symposium (RTSS)}, 1994.

\bibitem{spuri1996}
M.~Spuri and G.~Buttazzo,
``Scheduling Aperiodic Tasks in Dynamic Priority Systems,''
\textit{Real-Time Systems}, vol.~10, no.~2, March 1996.

\bibitem{lss1987}
J.~Lehoczky, L.~Sha, and J.~Strosnider,
``Enhanced Aperiodic Responsiveness in Hard Real-Time Environments,''
in \textit{Proc. IEEE Real-Time Systems Symposium (RTSS)}, 1987.

\bibitem{cc1989}
H.~Chetto and M.~Chetto,
``Some Results of the Earliest Deadline Scheduling Algorithm,''
\textit{IEEE Trans. Software Engineering}, vol.~15, no.~10, 1989.

\bibitem{ai_usage_m3}
S.~Boiddo,
\textit{AI Usage Protocol, Milestone 3: Improved Priority Exchange Server},
HSHL, 2026. Protocol version 0.2; JSON and BibTeX files submitted with this
milestone.

\end{thebibliography}

\end{document}
